\section{A common upper bound in the full parameter neighborhood}
\label{sec:parameter-cover}\label{sec:upper-main}
The preceding propositions bound all counted cycles of a fixed field
on open coefficient windows. To obtain Theorem~\ref{thm:main}(i), we
must cover every raw parameter direction and choose common phase,
parameter, and Hausdorff neighborhoods. The difficulty is that the
strict simple-contact margins vanish at the positive-quadratic and
two-endpoint directions. We cover these boundary directions first,
then take minima on the remaining compact set.

Use the radius $\varepsilon$ and direction $(E_0,E_1,E_2,D)\in S^3$
defined in Section~\ref{sec:return-framework}. The root-collision and
two-endpoint estimates are supplemented by the
exclusions and positive-value estimate of Appendix~\ref{sec:whole-period},
the simple-contact comparison of Appendix~\ref{sec:simple-contact},
and the forcing estimates of Appendix~\ref{sec:contact-forcing}.

\subsection{The positive quadratic direction}
Near $(E_0,E_1,E_2,D)=(0,0,1,0)$, a second weighted radius separates
the forcing, constant, and linear terms from the quadratic term.
Replacing
\[
 \varepsilon\mapsto\varepsilon\sqrt{E_2},\qquad
 (E_0,E_1)\mapsto(E_0,E_1)/E_2,\qquad D\mapsto D/E_2^{3/2}
\]
normalizes $E_2=1$ without changing the field. Write
\begin{equation}\label{eq:coefficient-sphere}
 R=\max(|E_1|,\sqrt{|E_0|},|D|^{1/3}),\quad
 P=z^2+Ra z+R^2c,\quad D=R^3d.
\end{equation}
At $R=0$ there is no positive periodic orbit by
Lemma~\ref{lem:damping}. Otherwise
$\max(|a|,\sqrt{|c|},|d|^{1/3})=1$. Use the six faces in
the order $d=1,d=-1,c=-1,c=1,a=1,a=-1$, assigning ties once.
Proposition~\ref{prop:scaled-positive-forcing} bounds the first face
by one, and Lemma~\ref{lem:scaled-negative-gates} excludes the next
two. Proposition~\ref{prop:positive-constant} treats $c=1$.
On $a=1$, $c\le1/8$
is excluded; for $c\ge1/8$ the substitution
\[
r=R\sqrt c,\qquad a_*=1/\sqrt c,\qquad d_*=d/c^{3/2}
\]
again gives the compact positive-constant family.

The exclusion on $a=1$ also holds on an open coefficient set.  For
$a>15/16$ and $c<5/32$, take $\alpha=1/2$ in
Lemma~\ref{lem:scaled-negative-gates}, since
\[
 R^{-2}P(-R/2)=\tfrac14-\tfrac a2+c<-\tfrac1{16}.
\]
For $c>1/16$ use the positive-constant substitution above, with $a$
in a fixed neighborhood of one.  These two open sets overlap and
cover the whole $a=1$ face.

These face estimates extend to open coefficient windows. For $d>0$ put
\[
 \widehat R=R d^{1/3},\qquad
 \widehat a=a/d^{1/3},\qquad
 \widehat c=c/d^{2/3},\qquad \widehat d=1 .
\]
On a neighborhood with $d$ bounded below, the hatted coefficients range
in a fixed compact set. Proposition~\ref{prop:scaled-positive-forcing}
therefore treats a neighborhood of the $d=1$ face, including its ties
with $a$ or $c$. If $c>0$ is bounded below, the substitution
\(\widehat R=R\sqrt c\), \(\widehat a=a/\sqrt c\), and
\(\widehat d=d/c^{3/2}\) gives the positive-constant family.  The strict
negative margins defining the $d=-1$ and $c=-1$ faces persist on their
respective neighborhoods. Ties are assigned after these overlapping
open regions have been chosen.

For a neighborhood of the remaining face $a=-1$, take $a<0$
bounded away from zero and put

\[
 A=-Ra>0,\qquad \widehat c=\frac{c}{a^2},\qquad
 \widehat d=\frac{d}{(-a)^3}.
\]

After the change $z=A\,\widehat z$, the polynomial and forcing become
$P=A^2(\widehat z^2-\widehat z+\widehat c)$ and
$D=A^3\widehat d$. On the face itself $A=R$.
Choose numbers
\(0<c_{\rm in}<c_{\rm out}\) and
\(0<d_{\rm in}<d_{\rm out}\), all sufficiently small, so that
Proposition~\ref{prop:simple-contact} applies on the inner box
$|\widehat c|<c_{\rm out}$, $|\widehat d|<d_{\rm out}$, with $C=1$.
The outer reductions start at the smaller thresholds $c_{\rm in}$
and $d_{\rm in}$, so their windows overlap this box.
If $|\widehat d|>d_{\rm in}$, change radius
to $A|\widehat d|^{1/3}$. The new forcing is $1$ or $-1$; the other
coefficients, $-|\widehat d|^{-1/3}$ and
$\widehat c|\widehat d|^{-2/3}$, range in fixed
compact sets, to which the forcing results apply. If
$|\widehat d|<d_{\rm out}$
and $\widehat c<-c_{\rm in}$, change radius to
$A\sqrt{|\widehat c|}$ and apply
the negative-constant exclusion. If $|\widehat d|<d_{\rm out}$ and
 $\widehat c>c_{\rm in}$,
put
\[
 r=A\sqrt{\widehat c},\qquad a_*=-1/\sqrt{\widehat c},\qquad
 d_*=\widehat d/\widehat c^{3/2},
\]
and use Proposition~\ref{prop:positive-constant}, including its
separated positive roots and their collision. Each new coefficient
range is compact and fixed before the corresponding radial scale is made
small. Each substitution reaches one of the stated regional estimates,
so the six-face decomposition ends here.

Each parameter therefore belongs to an open window whose estimate
bounds all its counted cycles by two. At a positive-root collision,
localization and finite-event gluing count the through and fiber-exit
pieces together on one source interval. Elsewhere the bound is no
larger. Fix the physical sections and coefficient thresholds, then
take the finitely many required shrinkings of $R$, $\varepsilon$,
the source interval, and the Hausdorff neighborhood. This gives an
open-neighborhood bound of two at the positive-quadratic direction,
with arbitrary relative rates of the coefficients.

\subsection{The remaining normalized directions}
Classify the other directions first by the signs of the polynomial
and forcing, then by the endpoint-root configuration.
Fix $b_*$ and let $P_*,D_*$ denote a direction on $S^3$.
Lemma~\ref{lem:negative-gates} excludes an open neighborhood
unless $P_*\ge0$ on $[f_{b_*},0]$ and $D_*\ge0$.
If $P_*(0)>0$, Proposition~\ref{prop:positive-value} gives a
neighborhood bound of one. If $P_*(0)=0$ and $D_*>0$,
Proposition~\ref{prop:positive-forcing} gives a bound of two.
The remaining directions satisfy
\begin{equation}\label{eq:remaining-polynomial}
 E_{0*}=D_*=0,\qquad
 P_*(z)=z(-A_*+C_*z),\qquad A_*^2+C_*^2=1.
\end{equation}
For $x=-z\in[0,-f_{b_*}]$,
$P_*(-x)=x(A_*+C_*x)\ge0$. Hence $A_*\ge0$. If $A_*=0$,
then $C_*=1$, the positive-quadratic direction treated above. If $A_*>0$ and
$P_*>0$ on $[f_{b_*},0)$, the negative preceding interval has
a fixed positive polynomial margin. In a neighborhood use
$A=-E_1$, $c=E_0/A^2$, $d=D/A^3$ in
Proposition~\ref{prop:simple-contact}; here $A$ is bounded away
from zero and $C=E_2$ may have either sign. This gives at most two.
The only other possibility is
\[
 A_*>0,\qquad C_*<0,\qquad C_*f_{b_*}=A_*.
\]
The roots $f_{b_*}$ and zero are simple, with derivatives $A_*$
and $-A_*$. Proposition~\ref{prop:two-endpoints} supplies a
neighborhood bound of two. This list is exhaustive. In particular,
an interior double root or a double root at the lower endpoint must
have $P_*(0)>0$ and has already been treated. If $P_*\equiv0$,
normalization forces $D_*=1$ or $-1$, also already treated.

\subsection{Necessary directions and the order of the residual cover}
\label{sec:effective-residual-cover}
Fix a compact interval $J\Subset(0,2)$.
Let $\mathcal R_J^-$ be the compact set of possible limiting
directions, including the modulus $b\in J$, of negative-forcing
whole-graphic cycles. The gate and damping identities give
\[
 D=0,\qquad E_1\le0,\qquad
 P(z)\ge0\quad(f_b\le z\le0).
\]
In the zero-constant subset $E_0=0$ these reduce to
$P=z(-A+Cz)$, $A\ge0$, $A^2+C^2=1$.
The only non-strict simple-contact boundaries are $A=0,C=1$
and $A>0,C<0,Cf_b=A$.

More precisely, for every relatively open
$N\supset\mathcal R_J^-$ there exist $\varepsilon_N,\rho_N>0$
such that every $D<0$ field with
$0<\varepsilon<\varepsilon_N$ and a cycle within
$\rho_N$ of the whole graphic has direction in $N$.
If not, choose $\varepsilon_n,\rho_n\to0$ and counterexamples
outside $N$. Compactness gives a direction subsequence
$(b_n,q_n)\to(b_*,q_*)$. A strictly negative limiting $D_*$ or
a negative value of $P_*$ is excluded by
Lemma~\ref{lem:negative-gates}; continuity of the positive
half-plane damping restriction gives $E_{1*}\le0$.
Their whole-Hausdorff convergence supplies the descending crossing
required by the gate argument.
Thus the limit belongs to $\mathcal R_J^-\subset N$,
contradicting the complement's closedness.
The two thresholds depend on the chosen $N$.

Choose effective smaller positive-quadratic and two-endpoint
windows first. Their complements in the compact zero-constant
necessary set form a compact strict simple-contact remainder.
Only on that remainder fix the positive minima of $A$ and
the negative preceding polynomial margin.
Finitely many strict simple-contact windows cover it.
The resulting open union contains the zero-constant necessary
set. Compactness then gives an $\eta>0$ such that the
remaining necessary directions satisfy $E_0\ge2\eta$.
Cover these by finitely many positive-value windows with fixed
$E_0$ margin. Now apply the preceding necessary-direction
statement to this effective open union.
Directions with $D\ge0$ are covered directly by the
negative-polynomial gate, positive-value, positive-forcing,
simple-contact, and two-endpoint windows.
Thus every strict minimum is taken on a fixed compact remainder
after the boundary windows have been chosen.

\begin{lemma}[Compactness of the graphic family]
\label{lem:compact-graphic-family}
For every compact $J\Subset(0,2)$, the set
\[
 \mathcal G_J:=\bigcup_{b\in J}\Gamma_0^b
\]
is compact.  More precisely, let $J_\nu\Subset(0,2)$ be finitely many
compact subintervals and, for each $\nu$, let
$\mathcal O_{\nu i}$ ($i\in\mathcal I_\nu$) be the finitely many local phase
neighbourhoods used by the coefficient charts over $J_\nu$.  If
$\mathcal G_{\nu i}$ denotes the union of the corresponding graphic
family and the finitely many section cores for chart $(\nu,i)$, and
$\mathcal G_{\nu i}\Subset\mathcal O_{\nu i}$, then
\[
 \delta_{\nu i}:=\operatorname{dist}
   (\mathcal G_{\nu i},\,\mathbb R^2\setminus\mathcal O_{\nu i})>0,
 \qquad
 \delta_*:=\min_{\nu,i}\delta_{\nu i}>0.
\]
The global compact phase neighborhood $K$ contains the union of the
graphic families and section cores. Each chart uses its own local
neighborhood $\mathcal O_{\nu i}$ and the common clearance $\delta_*$.
\end{lemma}
\begin{proof}
The explicit fast orbit in~\eqref{eq:fast-orbit} and the implicit-function
construction of $t_b$ and $f_b$ in Section~\ref{sec:geometry} make
$b\mapsto t_b$ and $b\mapsto f_b$ continuous on $J$.  Hence the two maps
\[
 (b,s)\longmapsto \gamma_b(st_b),\qquad
 (b,s)\longmapsto (sf_b,0),
 \qquad (b,s)\in J\times[0,1],
\]
are continuous and have compact images. Their union is exactly
$\mathcal G_J$.  For each fixed pair $(\nu,i)$, the same argument applies
to the compact set $\mathcal G_{\nu i}$, so its distance to
 $\mathbb R^2\setminus\mathcal O_{\nu i}$ is positive.  There are only finitely many
pairs, and hence their minimum $\delta_*$ is positive.
\end{proof}

\begin{proof}[Proof of assertion (i) of Theorem~\ref{thm:main}]
Fix $J\Subset(0,2)$. At each $b_0\in J$, take a finite subcover of
$S^3$ by the open coefficient windows above, intersect their modulus
intervals and take the minima of their radial and Hausdorff thresholds.
Choose each chart's phase neighborhood to contain the graphic and
closed enlarged section cores. Their continuity
(Lemmas~\ref{lem:compact-graphic-family} and~\ref{lem:sections})
allows a common shrinking of this modulus interval preserving these
inclusions over its closure. Now take a finite cover of $J$ by such
intervals $I_\nu$, and compact subintervals $J_\nu\Subset I_\nu$
still covering $J$. The corresponding finite chart sets
$\mathcal I_\nu$ and fixed phase neighborhoods satisfy
$\mathcal G_{\nu i}\Subset\mathcal O_{\nu i}$, with the notation of
Lemma~\ref{lem:compact-graphic-family}.

By Lemma~\ref{lem:compact-graphic-family}, choose a compact phase
neighborhood $K$ of $\mathcal G_J$ containing the finitely many required
transverse sections. The same lemma gives a common positive clearance
$\delta_*$ for the chartwise inclusions
$\mathcal G_{\nu i}\Subset\mathcal O_{\nu i}$. Choose the common
Hausdorff radius below $\delta_*$, all local phase clearances, and all
section-core radii. Let $\varepsilon_*$ be the minimum of the positive
radial thresholds, and take
\begin{equation}\label{eq:raw-neighborhood-main}
 V=\left\{\mu:
 \frac{\mu_1^2+\mu_2^2+\mu_3^2}{\varepsilon_*^4}
 +\frac{\mu_4^2}{\varepsilon_*^6}<1\right\}.
\end{equation}
Equation~\eqref{eq:weighted-radius-main} puts every nonzero $\mu\in V$ below
the common radial threshold.

For fixed $b\in J$ and nonzero raw parameter, select a compact core
$J_\nu$ containing $b$ and an applicable coefficient window. Its regional
result bounds the total actual multiplicity of all isolated cycles at the
chosen Hausdorff distance: Lemma~\ref{lem:sections} places them on the
same signed source interval. Every regional bound is at most two, so
selecting any applicable window gives the bound two for this field.
Each regional result already counts multiplicities on all its
transition intervals and finite seams;
in particular, a quadratic late zero cannot be accompanied by an early
or threshold zero. At $\mu=0$, Lemma~\ref{lem:damping} excludes positive
periodic orbits. Finite transition endpoints are included; equilibrium
limits and infinite-time passages are not periodic orbits. Thus one
$K,V,\rho_{\rm H}$ works for the whole interval $J$.
\end{proof}
